\section{Conclusion}
In this paper, we introduced the Sequence-Single Index model as a new, high-dimensional theoretical framework for analyzing single-layer attention architectures. The most significant contribution of this study is the definition of the Sequence Information Exponent. This exponent, which serves as a rigorous tool to quantify the inherent hardness and predict the sample complexity of Stochastic Gradient Descent, is defined in direct analogy with classical single-index models. This analysis transcends qualitative understanding, providing precise scaling laws for the required number of gradient steps, denoted by $n$, relative to the token dimension, denoted by $d$. We demonstrated that the sequential setting is significantly richer, showing how the sequence length $L$ accelerates convergence and how positional encoding can proactively reduce the SIE, thereby lowering the computational barrier to learning.

This work establishes a foundational line of research essential for a principled understanding of modern sequential data models, particularly those based on the Transformer architecture. The intricate dynamics manifesting in the population loss landscape, exemplified by the identification of phase transitions and the convergence to suboptimal local minima, unveil a wealth of avenues for future investigation. Subsequent research should aim to fully map these complex high-dimensional dynamics, develop techniques for robustly breaking inherent symmetries, and extend the SIE framework to more complex multi-index and multi-layer attention systems. We hope that this quantitative framework will stimulate further theoretical explorations, leading to the development of a robust, generalizable understanding of learning on structured sequential data.

\paragraph{Limitations ---} The theoretical analysis of SGD in Sequence Single-Index models is subject to several simplifications for tractability. Specifically, the model under scrutiny is a simplified, single-layer attention architecture where the key and query matrices are tied, the attention head dimension is one (\(d_\mathrm{head}=1\)), and the value matrix is the identity (\(V=I_L\)). It should be noted that the scope of this architecture is restricted to a specific class of target functions, including even functions, although Appendix~\ref{2506:app:sie-extra} provides an extension to this class. Additionally, the analysis operates under the assumption of token independence across the sequence, a crucial simplification that facilitates the execution of numerous sequence modeling tasks. The analysis of SGD dynamics is founded on the approximation of the discrete dynamics via a second-order ODE (Gradient Flow approximation).
